\section{Panel software as a contract for timing, weighting and uncertainty}
\label{exec23:panel-scope}

A bank comparing staggered branch interventions with a simultaneous regional
launch faces two different designs. A Python fixed-effect regression and an R
synthetic difference-in-differences fit can both return a scalar measured in
dollars, but that coincidence does not establish that they estimate the same
policy value. The reusable object must describe the population, adoption clock,
control eligibility, outcome horizon, weighting and uncertainty. This section
examines concrete interfaces in PyFixest and the original R synthdid package,
building on the separate R fixest discussion. Source inspection establishes
what the selected code says and how a proposed adapter must represent it; it
does not establish that an installed package has passed our trials.
\citep{exec23_pyfixest_reuse_source,exec23_synthdid_reuse_source}

The PyFixest snapshot is commit
\texttt{45d52b6e368cefc7321239921f2547b4e08c8f20}, dated 6 October 2026.
Its Python version is dynamic through Maturin; the Rust core declares 0.60.0.
The original synthdid snapshot is commit
\texttt{70c1ce3eac58e28c30b67435ca377bb48baa9b8a}, dated 15 January 2024,
and declares version 0.0.9. The different dates matter for maintenance planning:
the latter README describes a beta interface, while the former has changing
Python and native numerical boundaries. A cutoff-selected commit and a package
version are identities, not promises of future support. Neither Python syntax
similarity to R fixest nor the use of the name SDID by other implementations
transfers a qualified capability.
\citep{exec23_pyfixest_reuse_source,exec23_synthdid_reuse_source}

\subsection{A fixed-effect coefficient can reverse a policy conclusion}
\label{exec23:panel-projection}
Let $D_{it}$ indicate treatment and $Y_{it}$ be a monetary outcome for unit $i$
at period $t$. In a balanced panel, residualizing $D$ against unit and time
effects gives
\begin{equation}
 \widetilde D_{it}=D_{it}-\bar D_{i\cdot}-\bar D_{\cdot t}+\bar D,
 \qquad
 \widehat\beta_{\rm TWFE}=
 \frac{\sum_{it}\widetilde D_{it}Y_{it}}
      {\sum_{it}\widetilde D_{it}^{\,2}}.
 \label{exec23:panel-twfe}
\end{equation}
This is a projection identity. To interpret it as the intended average causal
effect requires additional restrictions on treatment effects, counterfactual
trends and comparisons. The fixed-effect labels in a formula do not supply
those restrictions.

An exact constructed example makes the distinction visible without estimating
anything in PyFixest. Stipulate zero untreated potential outcomes and use
\begin{equation}
 D=\begin{pmatrix}0&1&1\\0&0&1\end{pmatrix},
 \qquad
 Y=\begin{pmatrix}0&1&10\\0&0&1\end{pmatrix},
 \qquad
 \widetilde D=\frac16
 \begin{pmatrix}-1&2&-1\\1&-2&1\end{pmatrix}.
 \label{exec23:panel-witness}
\end{equation}
All three treated-cell effects are positive and their average is $4$. The
denominator in (\ref{exec23:panel-twfe}) is $1/3$, so the projection is $-7/2$.
Its weights on the six outcomes are
$(-1/2,1,-1/2;1/2,-1,1/2)$. In particular, the large positive late effect
receives a negative weight. The example intentionally has no never-treated
unit and no untreated comparison in its final period; it demonstrates the
projection's meaning, not a successful causal design. With a stipulated cost
of $2$ dollars per treated cell, the target net is $2$ dollars but treating the
projection as that target would report $-11/2$ dollars. These numbers reproduce
with rational arithmetic and are neither native package outputs nor measured
bank returns. \citep{exec23_pyfixest_reference}

In the inspected PyFixest event-study dispatcher, the concrete TWFE branch
constructs a treatment indicator, optional covariates and unit/time effects.
The stored ATT flag does not itself generate a lead--lag formula in that
branch. DID2S and the saturated implementation follow different fitting and
aggregation paths. The product therefore needs an explicit estimator profile
and a displayed estimand, rather than an event-study label that conceals the
formula actually fitted. \citep{exec23_pyfixest_reuse_source}

\subsection{Calendar conversion, first-stage support and aggregation}
\label{exec23:panel-clock}
PyFixest's common DID base accepts numeric time and cohort fields but constructs
relative time by subtraction. A calendar code is not necessarily a period
index: $202601-202512=89$, whereas the elapsed month count is one. The proposed
adapter should retain real calendar dates and an explicit consecutive index,
including a policy for missing observations. The base constructs absorbing
treatment from the declared cohort and uses zero for never treated. Its event
origin is the first observed treated period. If declared adoption is period
two and the first observed post-treatment record is period three, the declared
event time is one while recentering on the observed record makes it zero. Both
dates must survive conversion so that an observation gap cannot silently
change the scientific question.
\citep{exec23_pyfixest_reuse_source,exec23_pyfixest_reference}

The DID2S helper verifies two unique numeric treatment values and then casts
them to Boolean. Codes two and three both map to true; the adapter must
validate the actual binary convention before that conversion. Its first stage
fits untreated observations. Missing fixed-effect levels in that sample lead
to missing predictions and an explicit error. This is meaningful support
information: a bank should retain which branches or calendars could not be
represented, instead of deleting the error from the request denominator. The
GMM covariance uses residuals from both stages. The wrapper retains the
second-stage inference degrees of freedom and disables specified
post-estimation capabilities; availability of a generic fitted-object method
does not establish that it is valid for this two-stage result.
\citep{exec23_pyfixest_reuse_source}

The saturated branch also needs a narrower contract than its common public
surface might suggest. The inspected dispatcher does not pass the covariate
formula to its saturated helper. The actual aggregation implementation
accepts period and shares choices, despite broader options in its docstring;
its period-weight code includes relative period zero despite contrary prose.
Its linear-combination interval uses the normal multiplier $1.96$. This is
not automatically a few-cluster correction or a simultaneous event-time band.
Those source discrepancies are reasons to test and constrain the chosen
profile, not empirical claims that a newly executed native request failed.
\citep{exec23_pyfixest_reuse_source}

Aggregation is itself part of the estimand. For effects
$\widehat\theta=(10,-2)'$, weights $w=(1/4,3/4)'$ and the stated covariance
\begin{equation}
 V=\begin{pmatrix}4&1\\1&9\end{pmatrix},
 \qquad w'\widehat\theta=1,
 \qquad w'Vw=\frac{91}{16},
 \label{exec23:panel-aggregation}
\end{equation}
the equal-cohort effect would instead be $4$. Omitting cross-covariance gives
$85/16$, a different variance even with unchanged marginal standard errors.
The example assumes fixed aggregation weights. If the weights are estimated,
their additional uncertainty must be justified for the chosen procedure.
The reporting schema should retain cohort sizes, eligible populations, full
covariance and coefficient order so an analyst can see why the summaries
differ. \citep{exec23_pyfixest_reference}

PyFixest offers useful numerical machinery: its demeaning wrapper checks the
success flag and raises on nonconvergence. But a profile must also retain
singleton removal, analytic-weight defaults, rank tolerance, solver choice,
backend-specific preconditioner fallback and warnings. Rust, within, optional
Numba and Torch paths are not one interchangeable execution guarantee.
The earlier weighted-IV source discussion remains relevant to its IV branch;
fixed-effect fitting and first-stage diagnostics do not by themselves provide
weak-instrument confidence-set inversion.
\citep{exec23_pyfixest_reuse_source}

\subsection{The original synthetic difference-in-differences block}
\label{exec23:panel-sdid}
The original synthdid package addresses a different layout. There are $N_0$
control units followed by $N_1$ treated units, and $T_0$ pre-treatment periods
followed by $T_1$ post-treatment periods. Its treatment indicator equals one
only in the bottom-right block. This is the block design identified in the
original paper's inspected discussion; broader extensions in the literature
are not evidence that this package supports staggered adoption.
\citep{synthdid2021,exec23_synthdid_reuse_source}

Let $Z=Y-X\beta$ denote outcomes after the selected covariate adjustment, and
let $\omega$ and $\lambda$ be control-unit and pre-period weights. The code
computes the bilinear contrast
\begin{equation}
 \widehat\tau=
 (-\omega',N_1^{-1}\mathbf1_{N_1}')
 Z
 (-\lambda',T_1^{-1}\mathbf1_{T_1}')',
 \label{exec23:panel-contrast}
\end{equation}
so treated units and post periods receive equal weights within that block.
A branch-period effect is not automatically a customer-weighted portfolio
effect. If customer populations differ, the proposal must either defend this
target or derive and qualify an alternative instead of changing weights
after seeing the answer. The code's scalar result alone cannot establish
no anticipation, no interference, or the counterfactual trend, factor and
noise restrictions needed for a causal interpretation.
\citep{exec23_synthdid_reuse_source}

The panel converter rejects missing values, nonbinary treatment and anything
other than one observation at every unit--time pair. It then checks
simultaneous adoption. It converts factor and date fields to character before
sorting. ISO dates may preserve calendar order, but arbitrary period labels
need an explicit ordering map. The option that leaves treated units unsorted
does not change the estimator's expectation that controls occupy the first
$N_0$ rows. Nor does the converter retain other columns as a covariate array;
that alignment remains an adapter obligation. The direct matrix entrypoint
has a different, narrower set of checks. The proposed worker must enforce
finite data, nonempty treatment/control and pre/post groups, valid noise
estimation support, and row identity on either entry route.
\citep{exec23_synthdid_reuse_source}

For an independent arithmetic reference, use
\begin{equation}
 Y=\begin{pmatrix}
 10&12&14&16\\20&22&26&28\\15&17&25&27\\18&20&28&30
 \end{pmatrix},
 \quad\omega=(3/4,1/4)',\quad\lambda=(1/4,3/4)',
 \label{exec23:panel-block-example}
\end{equation}
with two controls, two treated units and two periods on either side. The
treated-minus-weighted-control gap is four dollars in both pre-periods and
$19/2$ dollars in both post-periods. Thus (\ref{exec23:panel-contrast}) is
$11/2$ dollars; a stipulated four-dollar cost leaves $3/2$ dollars. These are
supplied weights, not weights fitted by synthdid. Their effective donor count
$1/\sum_i\omega_i^2=8/5$ describes concentration, not the number of
independent observations or guaranteed inferential degrees of freedom.
\citep{exec23_synthdid_reference}

The same example supplies a useful conversion oracle. Valid unit-sum control
weights cancel a common 100-dollar post-period shock. Replacing them by
$(1/2,1/4)'$, which sums to $3/4$, changes the contrast by 25 dollars. The
direct estimator checks supplied weight lengths but does not enforce the
simplex constraints for fixed weights. An SDID profile must check those
constraints explicitly. The synthetic-control wrapper deliberately uses
zero pre-period weights and no unit-weight intercept, so its convention must
remain a separate profile rather than being rejected by an indiscriminate
SDID validator. \citep{exec23_synthdid_reuse_source,exec23_synthdid_reference}

\subsection{Optimization and inference are separate reproducibility obligations}
\label{exec23:panel-inference}
The default unit penalty scales with estimated pre-treatment first-difference
noise and $(N_1T_1)^{1/4}$; the time penalty uses a small multiplier. The
no-covariate path uses Frank--Wolfe steps with exact line search, followed by
optional sparsification and another optimization. Its stopping condition
uses objective decrease and an iteration limit. The covariate path alternates
weight updates with a $1/t$ gradient step and tests absolute objective change.
Neither returns a separate certified error bound just because iteration
stopped. A worker should retain objective histories, limits, finite-value
checks and an independent numerical comparison when feasible. The upstream
CVXR/ECOS reference provides an explicit simplex objective; the inspected
test compares one California example at tolerance $0.03$. That source test
is useful for planning, but has not been run here and cannot establish
general parity or calibrated uncertainty.
\citep{exec23_synthdid_reuse_source}

Reproducibility also extends beyond the fitted point estimate. The attached
options retain penalties, intercepts, update flags, minimum decrease and the
main iteration limit, but omit the sparsification function and the
pre-sparsification iteration limit. Resampling refits reconstruct calls from
these saved options. When a caller changed an omitted option, subsequent
refits can return to its default. The proposed adapter must record the full
fit configuration independently and replay it deliberately; preserving the
scalar and its attached setup alone is insufficient. This is a static
call-path observation requiring native regression tests.
\citep{exec23_synthdid_reuse_source}

The package offers bootstrap, jackknife and placebo variance estimators with
different support requirements. Bootstrap and jackknife return missing
uncertainty for one treated unit. The default fixed-weight jackknife also
returns missing uncertainty when only one control has nonzero weight; it
fixes the unit and time weights rather than reoptimizing them. Placebo
requires more controls than treated units and constructs artificial treated
groups from controls. That construction needs its own scientific
comparability argument. The package vignette explicitly cautions about the
single-treated placebo illustration. A default refusal must not become a
zero standard error, and a placebo standard deviation is not by itself a
randomization-test p-value. \citep{exec23_synthdid_reuse_source}

The bootstrap retries until it has the requested count of nonmissing
estimates, with no attempt limit in the inspected loop. For two controls and
two treated units, independent sampling of four rows with replacement draws
only controls or only treated units with probability
$2(1/2)^4=1/8$. Even if every mixed draw computes successfully, the expected
attempt count per acceptable composition is $8/7$. Further numerical missing
values can increase work or prevent termination; exceptions have a distinct
failure path. This calculation quantifies composition rejection only. It is
not a native timing result or a claim that every problematic numerical fit
returns missing rather than throwing an error. Record attempts, exclusions,
accepted draws and deadlines separately.
\citep{exec23_synthdid_reuse_source,exec23_synthdid_reference}

Finally, a display method can change the inference presented to a reader.
The default summary calls bootstrap, a fast summary calls jackknife, and
the formatter uses the fast summary with a $1.96$ multiplier. A service that
requests placebo variance and later parses the formatted object could thus
report a different uncertainty convention. The output should carry the
selected variance method, seed, successful and failed draws, weights and
interval convention as explicit fields. Native output, statistical validity
and user-interface presentation need separate checks.
\citep{exec23_synthdid_reuse_source}

\subsection{A proposed reuse decision with accountable costs}
\label{exec23:panel-reuse}
For this proposal, both packages are retained as reference implementations
with conditional routes to a Python or R worker. PyFixest is attractive for
its formula interface and native numerical backends; synthdid exposes a
compact original block estimator and a separate optimization reference.
Direct adoption of either package as an unrestricted causal service is
deferred. Porting now would duplicate numerical, inferential and maintenance
obligations before a measured bottleneck establishes the benefit. A
cohort-time DID estimator may be preferable for staggered adoption, while
SC, ordinary DID and SDID must be compared under matched estimands and
appropriate assumptions rather than a shared scalar output.
\citep{exec23_pyfixest_reuse_source,exec23_synthdid_reuse_source}

PyFixest declares MIT licensing; synthdid declares GPL version two or later
or the supplied three-clause BSD alternative. These source declarations and
the retained notice texts do not clear every distributed dependency or
deployment artifact. The scientific owner must approve the target, timing,
support and uncertainty; the adapter maintainer must own identity conversion,
exact options, native oracles and upgrade regressions; the operator must
account for budgets, failure recovery and distribution requirements.
\citep{exec23_pyfixest_reuse_source,exec23_synthdid_reuse_source}

The first useful public or synthetic trial should therefore compare the
complete request: eligible sample construction, estimation, uncertainty,
failed requests, explanation and reproducible cost accounting. Measure
preparation time, model fitting, covariance or resampling, rejected draws,
worker recovery and maintenance separately. The exact examples here are
input/output oracles for such a trial, not evidence that it has already
succeeded. A successful source-based reuse decision can justify the
prototype experiment while leaving native qualification, external review
and actual bank economics explicitly unresolved.
\citep{exec23_pyfixest_reference,exec23_synthdid_reference}
